\documentclass[pdftex, 12pt]{article}

\usepackage[pdftex]{graphicx}
\usepackage{listings}
\usepackage{fancyhdr}
\usepackage{lastpage}
\usepackage{fullpage}
\usepackage{color}
\usepackage{amsmath}
\usepackage[pdftex,bookmarks=true,colorlinks=true,linkcolor=blue]{hyperref}
\usepackage{subfig}

%this sets the line at the header
\setlength{\headheight}{15.2pt}

\newcommand{\HRule}{\rule{\linewidth}{0.5mm}}

% This is all for formatting and making the Table of Contents according to 
% spec. Don't play with it.
\makeatletter
\renewcommand\l@section[2]{%
  \ifnum \c@tocdepth >\z@
    \addpenalty\@secpenalty
    \addvspace{1.0em \@plus\p@}%
    \setlength\@tempdima{1.5em}%
    \begingroup
      \parindent \z@ \rightskip \@pnumwidth
      \parfillskip -\@pnumwidth
      \leavevmode \bfseries
      \advance\leftskip\@tempdima
      \hskip -\leftskip
      #1\nobreak\ 
      \leaders\hbox{$\m@th\mkern \@dotsep mu\hbox{.}\mkern \@dotsep mu$}
     \hfil \nobreak\hb@xt@\@pnumwidth{\hss #2}\par
    \endgroup
  \fi}
\makeatother



\begin{document}

%import the title page
\begin{titlepage}
	\begin{center}

		% Upper part of the page
		\textsc{\LARGE University of Nevada, Reno}\\[.5cm]
		\includegraphics[width=0.15\textwidth]{./logo.png}\\[.5cm]

		\textsc{\large CS 302 | Data Structures } \\[.5cm]

		% Title
		\HRule \\[0.4cm]
		{ \huge \bfseries Assignment \#3}\\[0.4cm]

		\HRule \\[1.5cm]

		% Author and supervisor
		\begin{minipage}{0.4\textwidth}
			\begin{flushleft} \large
				\emph{Students:}\\
				Joshua \textsc{Gleason}\\
				Josiah \textsc{Humphrey}
			\end{flushleft}
		\end{minipage}
		\begin{minipage}{0.4\textwidth}
			\begin{flushright} \large
				\emph{Instructor:} \\
				Dr. George \textsc{Bebis}
			\end{flushright}
		\end{minipage}

		\vfill

		% Bottom of the page
		{\large \today}

	\end{center}

\end{titlepage}


%headers, footers, and table of contents
\pagestyle{fancy}
\renewcommand{\sectionmark}[1]{\markright{\thesection}}
\rhead{Page \thepage\ of \pageref{LastPage}}
\lhead{}
\lfoot{CS 302 | Spring 2010}
\cfoot{}
\renewcommand{\footrulewidth}{0.4pt}

\tableofcontents

%\listoffigures
\newpage

\lhead{Joshua Gleason \& Josiah Humphrey}
\rhead{Page \thepage\ of \pageref{LastPage}}
\rfoot{Section\ \rightmark}
\cfoot{}
\lfoot{CS 302 | Spring 2010}
\renewcommand{\footrulewidth}{0.4pt}

\section{Introduction}

In this assignment, we looked at 3 different data structures that are very important for computer science. We examined
the usefulness and implementation of the heap, the priority queue, and the binary search tree. The program that we
designed to test these data structures were fairly trivial, but the challenge was implementing the assignment using the
data structures that we were supposed to implement. 

The first part of the assignment was to make a binary search tree that help users and their passwords. The binary search
tree was built using the name of the user as the node name. This allowed us to quickly search the list of users. This is
efficient when there are many users and when the program user needs to find a specific user quickly. The implementation
of the binary search tree we implemented was a template class that could handle any data structure. For this assignment,
we made our own data structure that included the name and the pass as string objects. This allowed us to do many useful
things such as compare. This was useful because it allowed the quick comparison of the string objects when building the
binary search tree. 

In the second part of the assignment, we were asked to implement an updated priority queue that allowed the users to
remove a specific item from the tree and to update the elements within the tree. This was done through inheritance. A
class called ``pqueue.h'' was defined that was a basic priority queue. The new updated queue, U\_PQType, was derived
from this class in order to include the two new functions. This derived class used the heap property of the priority
queue to implement the assignments challenge. We were asked to read in numbers from a file and place them into a heap.

\section{Use of Code}

\section{Functions}

%%%%%%%%%%%%%%%%%%%%%%%%%%%%%%%%%%%%%%%%
%%%%%%%       PART 1      %%%%%%%%%%%%%%
%%%%%%%%%%%%%%%%%%%%%%%%%%%%%%%%%%%%%%%%
\subsection{binaryTree.h}
\begin{description}

	\item{\textsc{binaryTree}}
		\begin{lstlisting}
			binaryTree();
		\end{lstlisting}
		\begin{description}

			\item{Purpose}

				Constructor, simply initializes root to NULL

			\item{Input}

				N/A

			\item{Output}

				N/A

			\item{Assumptions}

				N/A

		\end{description}
	\item{\textsc{~binaryTree}}
		\begin{lstlisting}
			~binaryTree();
		\end{lstlisting}
		\begin{description}

			\item{Purpose}

				Free all the memory in the tree

			\item{Input}

				N/A

			\item{Output}

				N/A

			\item{Assumptions}

				N/A

		\end{description}
	\item{\textsc{makeEmpty}}
		\begin{lstlisting}
			void makeEmpty();
		\end{lstlisting}
		\begin{description}

			\item{Purpose}
				
				Recursively remove all the nodes below the node.
			
			\item{Input}
			
				A node.

			\item{Output}

				N/A

			\item{Assumptions}

				Node is either NULL or points to an allocated node type.

		\end{description}
	\item{\textsc{isEmpty}}
		\begin{lstlisting}
			bool isEmpty() const;
		\end{lstlisting}
		\begin{description}

			\item{Purpose}

				Test if tree is empty.

			\item{Input}

				N/A

			\item{Output}

				Returns true if tree is empty, otherwise returns false.

			\item{Assumptions}

				N/A

		\end{description}
	\item{\textsc{numberOfNodes}}
		\begin{lstlisting}
			int numberOfNodes() const;
		\end{lstlisting}
		\begin{description}

			\item{Purpose}

				Count the number of nodes in the tree recursively.

			\item{Input}

				N/A

			\item{Output}

				Return the number of nodes in the tree.

			\item{Assumptions}

				N/A

		\end{description}
	\item{\textsc{retrieveItem}}
		\begin{lstlisting}
			bool retrieveItem(iType&);
		\end{lstlisting}
		\begin{description}

			\item{Purpose}

				Attempts to retrieve an item from the tree, return false if item is not found in the tree.

			\item{Input}

				An item that has the same "key" value as the item in question.

			\item{Output}

				Return true if item is in tree, false otherwise.

			\item{Assumptions}

				N/A

		\end{description}
	\item{\textsc{insertItem}}
		\begin{lstlisting}
			void insertItem(iType);
		\end{lstlisting}
		\begin{description}

			\item{Purpose}

				Create a new node and insert it into the tree with the value of the parameter.

			\item{Input}

				The value to be placed in the tree.

			\item{Output}

				N/A

			\item{Assumptions}

				N/A

		\end{description}
	\item{\textsc{deleteItem}}
		\begin{lstlisting}
			void deleteItem(iType);
		\end{lstlisting}
		\begin{description}

			\item{Purpose}

				Remove a node from the tree.

			\item{Input}

				An item with the same key value as the item to be removed.

			\item{Output}

				N/A

			\item{Assumptions}

				The value of the item must exist in the tree, use retrieveItem to test before calling this unless it is
				certain that the item is in the tree.

		\end{description}
	\item{\textsc{resetTree}}
		\begin{lstlisting}
			void resetTree(oType);
		\end{lstlisting}
		\begin{description}

			\item{Purpose}

				Generate a queue to traverse the tree in a particular order.

			\item{Input}

				An order type which can be IN\_ORDER, PRE\_ORDER, or POST\_ORDER, each one initializes a different
				queue.

			\item{Output}

				N/A

			\item{Assumptions}

				N/A

		\end{description}
	\item{\textsc{getNextItem}}
		\begin{lstlisting}
			bool getNextItem(iType&, oType);
		\end{lstlisting}
		\begin{description}

			\item{Purpose}

				Retrieve the next item in the queue corresponding to the value of oType.

			\item{Input}

				A reference parameter used to store the retrieved data and an order type of which queue to retrieve
				from.

			\item{Output}

				Return true if the item is the last item in the queue, also make the reference parameter equal to the
				value of the next item in the queue.

			\item{Assumptions}

				Assumes the proper queue as been initialized and is not empty.  Use the resetTree and isEmpty functions
				to ensure this is the case.

		\end{description}
	\item{\textsc{printTree}}
		\begin{lstlisting}
			void printTree(ostream&) const;
		\end{lstlisting}
		\begin{description}

			\item{Purpose}

				Prints the tree to the desired ostream type.

			\item{Input}

				An ostream object by reference, this could be cout or an ofstream.

			\item{Output}

				Outputs all values of the tree to the desired ostream.

			\item{Assumptions}

				Assumes that the left shift opperator << has been overloaded for ostream to output the iType values
				properly.

		\end{description}
	\item{\textsc{countNodes}}
		\begin{lstlisting}
			int countNodes(treeNode<iType>*);
		\end{lstlisting}
		\begin{description}

			\item{Purpose}

				The client function called by numberOfNodes that does the recursive counting.

			\item{Input}

				A node pointer.

			\item{Output}

				Recursively determines the number of nodes connected to the parameter node plus 1.

			\item{Assumptions}

				Assumes parameter points to valid node or is NULL.

		\end{description}
	\item{\textsc{retrieve}}
		\begin{lstlisting}
			bool retrieve(treeNode<iType>*, iType&);
		\end{lstlisting}
		\begin{description}

			\item{Purpose}

				The recursive component to retrieveItem.

			\item{Input}

				A node pointer and the item in question.

			\item{Output}

				Returns true if item is found, false if it is not.

			\item{Assumptions}

				Assumes parameter points to valid node or is NULL.

		\end{description}
	\item{\textsc{insert}}
		\begin{lstlisting}
			void insert(treeNode<iType>*&, iType);
		\end{lstlisting}
		\begin{description}

			\item{Purpose}

				The recursive component to insertItem, recursively traverses the tree to find the appropriate place to
				put the node.

			\item{Input}

				A node pointer and the item to be inserted.
				
			\item{Output}

				N/A

			\item{Assumptions}

				Assumes the pointer points to a valid node or is NULL.

		\end{description}
	\item{\textsc{deleteOut}}
		\begin{lstlisting}
			void deleteOut(treeNode<iType>*&, iType);
		\end{lstlisting}
		\begin{description}

			\item{Purpose}

				This portion of the function finds the correct node to be deleted recursively, then calls deleteNode to
				actually do the delete.

			\item{Input}

				A node pointer and the item to be deleted.

			\item{Output}

				N/A

			\item{Assumptions}

				Assumes the value held by the iType parameter exists in the tree.

		\end{description}
	\item{\textsc{deleteNode}}
		\begin{lstlisting}
			void deleteNode(treeNode<iType>*&);
		\end{lstlisting}
		\begin{description}

			\item{Purpose}

				Delete the node without breaking any rules of the binary search tree.

			\item{Input}

				The node the be deleted.

			\item{Output}

				N/A

			\item{Assumptions}

				Assumes the node is part of a valid binary search tree.

		\end{description}
	\item{\textsc{getPredecessor}}
		\begin{lstlisting}
			void getPredecessor(treeNode<iType>*, iType&);
		\end{lstlisting}
		\begin{description}

			\item{Purpose}

				Return the right most node of a left child.

			\item{Input}

				The left child of some node.

			\item{Output}

				Places the info from the right most child into the iType parameter.

			\item{Assumptions}

				Assumes the node is a valid node.

		\end{description}
	\item{\textsc{print}}
		\begin{lstlisting}
			void print(treeNode<iType>*, ostream&);
		\end{lstlisting}
		\begin{description}

			\item{Purpose}
				
				The recursive part of the print function, prints the tree inorder by printing the left child, the head,
				and then the right child.

			\item{Input}

				A node to be printed and an ostream object to print to.

			\item{Output}

				Prints all the values below the given node as well as the node using recursion.

			\item{Assumptions}
			
				Assumes the same as printTree as well as assuming the node is a valid node or NULL.

		\end{description}
	\item{\textsc{destroy}}
		\begin{lstlisting}
			void destroy(treeNode<iType>*&);
		\end{lstlisting}
		\begin{description}

			\item{Purpose}

				Recursively deallocate all values in the tree.

			\item{Input}

				A node pointer.

			\item{Output}

				N/A

			\item{Assumptions}

				Assumes the node pointer points to a valid node or is NULL.

		\end{description}
	\item{\textsc{copyTree}}
		\begin{lstlisting}
			void copyTree(treeNode<iType>*&, treeNode<iType>*);
		\end{lstlisting}
		\begin{description}

			\item{Purpose}

			\item{Input}

			\item{Output}

			\item{Assumptions}

		\end{description}
	\item{\textsc{preOrder}}
		\begin{lstlisting}
			void preOrder(treeNode<iType>*&, queue<iType>&);
		\end{lstlisting}
		\begin{description}

			\item{Purpose}
				
				Initialize the given queue with the values of the given tree traversed using pre-order traversal.

			\item{Input}

				A valid node pointer and an empty queue.

			\item{Output}

				Set the values of the queue to the pre-order traversal of the tree.

			\item{Assumptions}

				Assumes the queue is empty and the node pointer points to a valid node or is NULL.

		\end{description}
	\item{\textsc{inOrder}}
		\begin{lstlisting}
			void inOrder(treeNode<iType>*&, queue<iType>&);
		\end{lstlisting}
		\begin{description}

		\item{Purpose}
			
			Initialize the given queue with the values of the given tree traversed using in-order traversal.

		\item{Input}

			A valid node pointer and an empty queue.

		\item{Output}

			Set the values of the queue to the in-order traversal of the tree.

		\item{Assumptions}

			Assumes the queue is empty and the node pointer points to a valid node or is NULL.

		\end{description}
	\item{\textsc{postOrder}}
		\begin{lstlisting}
			void postOrder(treeNode<iType>*&, queue<iType>&);
		\end{lstlisting}
		\begin{description}

		\item{Purpose}
			
			Initialize the given queue with the values of the given tree traversed using post-order traversal.

		\item{Input}

			A valid node pointer and an empty queue.

		\item{Output}

			Set the values of the queue to the post-order traversal of the tree.

		\item{Assumptions}

			Assumes the queue is empty and the node pointer points to a valid node or is NULL.

		\end{description}

\end{description}
\subsection{user.h}
\begin{description}
	\item{\textsc{getName}}
		\begin{lstlisting}
			string getName() const
		\end{lstlisting}
		\begin{description}

			\item{Purpose}

			\item{Input}

			\item{Output}

			\item{Assumptions}

		\end{description}
	\item{\textsc{getPass}}
		\begin{lstlisting}
			string getPass() const
		\end{lstlisting}
		\begin{description}

			\item{Purpose}

			\item{Input}

			\item{Output}

			\item{Assumptions}

		\end{description}
	\item{\textsc{setName}}
		\begin{lstlisting}
			void setName( string& rhs )
		\end{lstlisting}
		\begin{description}

			\item{Purpose}

			\item{Input}

			\item{Output}

			\item{Assumptions}

		\end{description}
	\item{\textsc{setPass}}
		\begin{lstlisting}
			void setPass( string& rhs )
		\end{lstlisting}
		\begin{description}

			\item{Purpose}

			\item{Input}

			\item{Output}

			\item{Assumptions}

		\end{description}

	\item{\textsc{operator>}}
		\begin{lstlisting}
			bool operator>( const user& rhs )
		\end{lstlisting}
		\begin{description}

			\item{Purpose}

			\item{Input}

			\item{Output}

			\item{Assumptions}

		\end{description}
	\item{\textsc{operator<}}
		\begin{lstlisting}
			bool operator<( const user& rhs )
		\end{lstlisting}
		\begin{description}

			\item{Purpose}

			\item{Input}

			\item{Output}

			\item{Assumptions}

		\end{description}
	\item{\textsc{operator>=}}
		\begin{lstlisting}
			bool operator>=( const user& rhs )
		\end{lstlisting}
		\begin{description}

			\item{Purpose}

			\item{Input}

			\item{Output}

			\item{Assumptions}

		\end{description}
	\item{\textsc{operator<=}}
		\begin{lstlisting}
			bool operator<=( const user& rhs )
		\end{lstlisting}
		\begin{description}

			\item{Purpose}

			\item{Input}

			\item{Output}

			\item{Assumptions}

		\end{description}
	\item{\textsc{operator==}}
		\begin{lstlisting}
			bool operator==( const user& rhs )
		\end{lstlisting}
		\begin{description}

			\item{Purpose}

			\item{Input}

			\item{Output}

			\item{Assumptions}

		\end{description}

\end{description}
\subsection{part1.cpp}
\begin{description}

	\item{\textsc{readFile}}
		\begin{lstlisting}
			bool readFile( string fileName, binaryTree<user>& tree );
		\end{lstlisting}
		\begin{description}

			\item{Purpose}

			\item{Input}

			\item{Output}

			\item{Assumptions}

		\end{description}
	\item{\textsc{storeTree}}
		\begin{lstlisting}
			void storeTree( binaryTree<user>& tree, oType order, string fileName );
		\end{lstlisting}
		\begin{description}

			\item{Purpose}

			\item{Input}

			\item{Output}

			\item{Assumptions}

		\end{description}
	\item{\textsc{promptForMenu}}
		\begin{lstlisting}
			menuChoice promptForMenu();
		\end{lstlisting}
		\begin{description}

			\item{Purpose}

			\item{Input}

			\item{Output}

			\item{Assumptions}

		\end{description}
\end{description}
%%%%%%%%%%%%%%%%%%%%%%%%%%%%%%%%%%%%%%%%
%%%%%%%       PART 2      %%%%%%%%%%%%%%
%%%%%%%%%%%%%%%%%%%%%%%%%%%%%%%%%%%%%%%%
\subsection{heap.h}
\begin{description}

	\item{\textsc{reheapDown}}
		\begin{lstlisting}
			void reheapDown(int root, int bottom);
		\end{lstlisting}
		\begin{description}

			\item{Purpose}

			\item{Input}

			\item{Output}

			\item{Assumptions}

		\end{description}
	\item{\textsc{reheapUp}}
		\begin{lstlisting}
			void reheapUp(int root, int bottom);
		\end{lstlisting}
		\begin{description}

			\item{Purpose}

			\item{Input}

			\item{Output}

			\item{Assumptions}

		\end{description}
	\item{\textsc{swap}}
		\begin{lstlisting}
			void swap(ItemType &a, ItemType &b);
		\end{lstlisting}
		\begin{description}

			\item{Purpose}

			\item{Input}

			\item{Output}

			\item{Assumptions}

		\end{description}
\end{description}
\subsection{pqueue.h}
\begin{description}

	\item{\textsc{makeEmpty}}
		\begin{lstlisting}
			void makeEmpty();
		\end{lstlisting}
		\begin{description}

			\item{Purpose}

			\item{Input}

			\item{Output}

			\item{Assumptions}

		\end{description}
	\item{\textsc{isEmpty}}
		\begin{lstlisting}
			bool isEmpty() const;
		\end{lstlisting}
		\begin{description}

			\item{Purpose}

			\item{Input}

			\item{Output}

			\item{Assumptions}

		\end{description}
	\item{\textsc{isFull}}
		\begin{lstlisting}
			bool isFull() const;
		\end{lstlisting}
		\begin{description}

			\item{Purpose}

			\item{Input}

			\item{Output}

			\item{Assumptions}

		\end{description}
	\item{\textsc{enqueue}}
		\begin{lstlisting}
			void enqueue(ItemType newItem);
		\end{lstlisting}
		\begin{description}

			\item{Purpose}

			\item{Input}

			\item{Output}

			\item{Assumptions}

		\end{description}
	\item{\textsc{dequeue}}
		\begin{lstlisting}
			void dequeue(ItemType& item);
		\end{lstlisting}
		\begin{description}

			\item{Purpose}

			\item{Input}

			\item{Output}

			\item{Assumptions}

		\end{description}
\end{description}
\subsection{U\_PQType.h}
\begin{description}

	\item{\textsc{U\_PQType}}
		\begin{lstlisting}
			U\_PQType(int);
		\end{lstlisting}
		\begin{description}

			\item{Purpose}

			\item{Input}

			\item{Output}

			\item{Assumptions}

		\end{description}
	\item{\textsc{Remove}}
		\begin{lstlisting}
			void Remove(ItemType);
		\end{lstlisting}
		\begin{description}

			\item{Purpose}

			\item{Input}

			\item{Output}

			\item{Assumptions}

		\end{description}
	\item{\textsc{Update}}
		\begin{lstlisting}
			void Update(ItemType, ItemType);
		\end{lstlisting}
		\begin{description}

			\item{Purpose}

			\item{Input}

			\item{Output}

			\item{Assumptions}

		\end{description}
	\item{\textsc{printTree}}
		\begin{lstlisting}
			void printTree(std::ostream&);
		\end{lstlisting}
		\begin{description}

			\item{Purpose}

			\item{Input}

			\item{Output}

			\item{Assumptions}

		\end{description}
\end{description}
\subsection{part2.cpp}
\begin{description}
	\item{\textsc{readFile}}
		\begin{lstlisting}
			bool readFile( string fileName, U\_PQType<int>* &tree);
		\end{lstlisting}
		\begin{description}

			\item{Purpose}

			\item{Input}

			\item{Output}

			\item{Assumptions}

		\end{description}
	\item{\textsc{promptForMenu}}
		\begin{lstlisting}
			menuChoice promptForMenu();
		\end{lstlisting}
		\begin{description}

			\item{Purpose}

			\item{Input}

			\item{Output}

			\item{Assumptions}

		\end{description}

\end{description}
\section{Bugs and Errors}

\section{What was Learned}

\section{Division of Labor}

\section{Extra Credit}

\end{document}
